\documentclass[10pt,a4paper]{amsart}
\usepackage[margin=19mm]{geometry}
\usepackage{amsmath,amssymb,mathtools}
\usepackage[scr=rsfs,cal=euler]{mathalfa}
\usepackage{xfrac}
\usepackage[unicode,hidelinks,bookmarks=false]{hyperref}
\newcommand{\ie}{i.e.\ }
\newcommand{\cf}{cf.\ }
\newcommand{\resp}{respectively\ }
\newcommand{\Subsection}{\subsection}
\providecommand{\nobreakdash}{\nobreak\hbox{-}}
\setlength{\emergencystretch}{2em}
\multlinegap0pt
\usepackage{amssymb}
\usepackage{mathtools}
\newtheorem{lemma}{Lemma}
\newtheorem{proposition}{Proposition}
\newcommand{\one}{\mathbf{1}}
\newcommand{\rev}{\mathrm{rev}}
\title{A Linear Lower Bound for Dominating Sets in $k$-Majority Tournaments}
\author{Jiangdong Ai, Xiangjie Yi}
\date{Source: arXiv:2607.23148v1 (CC BY 4.0)}
\begin{document}
\maketitle

\noindent\emph{Source and licence.} A Linear Lower Bound for Dominating Sets in $k$-Majority Tournaments. Version arXiv:2607.23148v1 (CC BY 4.0), distributed by arXiv under the Creative Commons Attribution 4.0 International licence, http://creativecommons.org/licenses/by/4.0/, as recorded in the arXiv licence metadata for this exact version and shown on the abstract page https://arxiv.org/abs/2607.23148v1. Full licence notice: \texttt{licenses/arxiv-2607.23148-CC-BY-4.0.txt}.\par\smallskip

\noindent\textit{Editorial scope.} Counted unit: the article's proposition prop:four-orders (source0.tex lines 268--271). The article's complete original proof of it is delivered with it (source0.tex lines 273--339, one \texttt{proof} environment of 67 lines). No mathematics is written, repaired or re-derived: the statement and the proof are the article's own lines, in the article's own notation, and no retained line is substituted or re-flowed.  Retained uncounted as the source-local context the proof needs: lines 245--257.  Nothing was omitted from any retained span, so the package carries no omission record.  No retained line contains a citation, so the wrapper carries no bibliography and no bibliography entry.  No retained line contains a figure environment, an image-inclusion command or drawing code, so no image is delivered and the package declares no asset dependency.  Editorial formatting only, in addition: an AMS \texttt{amsart} preamble that restates the article's own declarations and macros used by the retained lines, namely the environments \texttt{lemma}, \texttt{proposition} on separate counters, together with the standard packages those lines need.  The retained environments print in this document as Lemma 1, Proposition 1 in order of appearance, whereas the article prints the same items with the numbering of its own full document.  The complete native source of the article as distributed in the arXiv e-print for this exact version is bundled unchanged as \texttt{sources/206006/source0.tex}.
\begin{lemma}\label{lem:psi-values}
For all admissible indices,
\[
  \psi(H_{r,a},H_{s,b})=0,
  \qquad
  \psi(V_{c,d},V_{e,f})=0,
\]
and
\begin{equation}\label{eq:cross-value}
  \psi(H_{r,a},V_{c,d})=
  \one_{\{a=c\}}-\one_{\{r=d\}}.
\end{equation}
\end{lemma}

\begin{proposition}\label{prop:four-orders}
The margin function $2\psi$ on $Z_q$ is realized exactly by four linear
orders.
\end{proposition}

\begin{proof}
For each $a\in[q]$, form the ordered blocks
\[
  \mathbf H^a=(H_{1,a},\ldots,H_{q,a}),
  \qquad
  \mathbf V^a=(V_{a,1},\ldots,V_{a,q}).
\]
We write blocks in succession to mean that every element of an earlier block lies above every element of a later one. Define the two column orders
\begin{align*}
  C_1&=\mathbf H^1\mathbf V^1\,
       \mathbf H^2\mathbf V^2\cdots
       \mathbf H^q\mathbf V^q,\\
  C_2&=(\mathbf H^q)^{\rev}(\mathbf V^q)^{\rev}\cdots
       (\mathbf H^1)^{\rev}(\mathbf V^1)^{\rev}.
\end{align*}
Note that $C_2$ is not the reverse of $C_1$: the pairs $\mathbf H^a\mathbf V^a$ occur in the reverse order and each block is
reversed internally, but $\mathbf H^a$ still precedes $\mathbf V^a$ for every $a$.

Consider first two states of the same type, say $H_{r,a}$ and
$H_{s,b}$.  If $a=b$, both lie in the block $\mathbf H^a$, whose
internal order is reversed in $C_2$; if $a\ne b$, the blocks
$\mathbf H^a$ and $\mathbf H^b$ occur in opposite orders in $C_1$ and $C_2$.  Either way the two contributions cancel, and the same argument
applies to vertical--vertical pairs. Hence the column pair has zero
margin on pairs of the same type.

Now consider $H_{r,a}$ and $V_{c,d}$.  If $a=c$, then $H_{r,a}$ lies above $V_{c,d}$ in both column orders, by the remark above.  If
$a\ne c$, the blocks $\mathbf H^a$ and $\mathbf V^c$ occur in opposite orders in $C_1$ and $C_2$.  Therefore
\begin{equation}\label{eq:column-pair}
  \sigma_{C_1}(H_{r,a},V_{c,d})+
  \sigma_{C_2}(H_{r,a},V_{c,d})
  =2\one_{\{a=c\}}.
\end{equation}

For each $r\in[q]$, form the blocks
\[
  \widehat{\mathbf V}^{\,r}=(V_{1,r},\ldots,V_{q,r}),
  \qquad
  \widehat{\mathbf H}^{\,r}=(H_{r,1},\ldots,H_{r,q}),
\]
and define the two row orders
\begin{align*}
  R_1&=\widehat{\mathbf V}^{\,1}\widehat{\mathbf H}^{\,1}\,
       \widehat{\mathbf V}^{\,2}\widehat{\mathbf H}^{\,2}\cdots
       \widehat{\mathbf V}^{\,q}\widehat{\mathbf H}^{\,q},\\
  R_2&=(\widehat{\mathbf V}^{\,q})^{\rev}
       (\widehat{\mathbf H}^{\,q})^{\rev}\cdots
       (\widehat{\mathbf V}^{\,1})^{\rev}
       (\widehat{\mathbf H}^{\,1})^{\rev}.
\end{align*}
Again $R_2$ is not the reverse of $R_1$: the pairs
$\widehat{\mathbf V}^{\,r}\widehat{\mathbf H}^{\,r}$ occur in the
reverse order and each block is reversed internally, but
$\widehat{\mathbf V}^{\,r}$ still precedes $\widehat{\mathbf H}^{\,r}$
for every $r$.  Exactly as above, the row pair has zero margin on pairs
of the same type.  For a cross pair, if $r=d$ then $V_{c,d}$ lies above
$H_{r,a}$ in both row orders, while if $r\ne d$ the blocks
$\widehat{\mathbf H}^{\,r}$ and $\widehat{\mathbf V}^{\,d}$ occur in
opposite orders.  Therefore
\begin{equation}\label{eq:row-pair}
  \sigma_{R_1}(H_{r,a},V_{c,d})+
  \sigma_{R_2}(H_{r,a},V_{c,d})
  =-2\one_{\{r=d\}}.
\end{equation}
Equations~\eqref{eq:column-pair} and~\eqref{eq:row-pair}, together with
Lemma~\ref{lem:psi-values}, show that $C_1,C_2,R_1,R_2$ have total
margin $2\psi$ on every pair of distinct states.
\end{proof}

\end{document}
